\section{Source-grounded atom-gravimeter binding}
\label{sec:apparatus}

\subsection{Baseline apparatus and acceleration response}

To convert the abstract resource variable into physical acquisition time, we use the compact three-pulse $^{87}$Rb Raman gravimeter reported by Le Gou\"et \emph{et al.} \cite{LeGouet2008} as the baseline apparatus family. The source reports an interferometer duration $2T=100$ ms, repetition rate $f_c=4$ Hz, best short-term phase noise $11$ mrad/shot, and acceleration sensitivity $1.4\times10^{-8}g$ at $1$ s. The compatible sensitivity-function formalism of Cheinet \emph{et al.} \cite{Cheinet2008} supplies the three-pulse transfer geometry.

At low frequency, the acceleration phase is
\begin{equation}
\Phi_a=k_{\rm eff}aT^2,
\label{eq:ai-scale}
\end{equation}
and the normalized response to sinusoidal acceleration at frequency $f$ is represented by
\begin{equation}
R_a(f)=\left[\frac{\sin(\pi fT)}{\pi fT}\right]^2.
\label{eq:ra}
\end{equation}
Using $\lambda=780$ nm, counter-propagating $k_{\rm eff}=4\pi/\lambda$, and $T=50$ ms gives
\begin{align}
k_{\rm eff}&=1.61107316\times10^7\ {\rm m}^{-1},\\
k_{\rm eff}T^2&=4.02768289\times10^4\ {\rm rad/(m\,s^{-2})}.
\end{align}
At four shots per second, the measured $11$ mrad/shot therefore implies $1.39247\times10^{-8}g$ at $1$ s, reproducing the published $1.4\times10^{-8}g$ scale to rounding accuracy. This reproduction is a dimensional consistency check, not a fit.

The science parameter $\theta$ is an acceleration-like RQIR amplitude. We intentionally modulate it in a declared analysis band. The local phase model also contains a common multiplicative phase/acceleration scale, constant, linear, and quadratic phase terms, and a reference amplitude when that reference is not externally fixed. Static acceleration-like science is retained as a required negative control.

\subsection{From a published vibration spectrum to sampled covariance}

The same apparatus paper reports the measured passive-platform acceleration spectrum and the propagation of vibration noise through the interferometer \cite{LeGouet2008}. In the repository calculation, a coarse log--log digitization of the passive-platform curve is frozen explicitly as an approximate source-traceable digitization, not as raw campaign data. No amplitude normalization is fitted to the integrated sensitivity.

For one-sided acceleration spectral density $S_a(f)$, the acceleration response and scale factor define a phase spectrum $S_\Phi(f)$. The sampled shot covariance at lag $\ell$ is then formed as
\begin{equation}
C_\ell=\int_0^{\infty}S_\Phi(f)
\cos(2\pi f\ell T_c)\,df,
\qquad T_c=f_c^{-1}.
\label{eq:covariance}
\end{equation}
The covariance is therefore derived at the actual shot spacing rather than replaced by a single arbitrary correlation coefficient.

Without tuning its amplitude, the digitized spectrum gives
\begin{equation}
\sigma_{a,\rm vib}(1\,{
m s})=6.34856\times10^{-8}g,
\end{equation}
compared with the source's inferred $6.5\times10^{-8}g$ vibration limit, a $2.33\%$ difference. Applying the reported factor-three vibration rejection and adding the independently evaluated $4$ mrad/shot non-vibration budget gives
\begin{equation}
\sigma_{a,\rm corr}(1\,{
m s})=2.17592\times10^{-8}g,
\end{equation}
compared with the reported typical $2\times10^{-8}g$ corrected scale, an $8.80\%$ difference. The agreement of two independent integrated scales, without re-normalizing the digitized spectrum, is used as the principal validation of the covariance surrogate.

\subsection{Physical colored-noise identifiability}

The source-traceable Toeplitz covariance is inserted directly into the joint Fisher/nuisance calculation. The same structural pattern found in the abstract analysis survives:
\begin{itemize}
\item static science plus a calibrated reference fails because the science direction is absorbed by a free phase offset;
\item modulated science without a calibrated absolute reference fails because of common multiplicative-scale degeneracy;
\item modulated science plus a known or finitely calibrated reference passes;
\item an unknown, unconstrained reference amplitude fails because it merely transfers the scale degeneracy into the science/reference ratio.
\end{itemize}

For the nominal physical spectrum, the phase-level forecast with finite reference calibration gives $\sigma_\theta=5.72247$ ng at $60$ s, $1.85019$ ng at $600$ s, and $0.804473$ ng at $3600$ s. The 2008 source also identifies a broad residual structure near $2$ Hz associated with horizontal-to-vertical crosstalk. To avoid a favorable-spectrum result, the repository multiplies that structure by smooth stress factors up to five. At $60$ s, the science uncertainty changes from $5.72$ ng at the nominal spectrum to $9.31$ ng for the five-fold stress while local science estimability remains nonzero.

These are design forecasts using published apparatus scales and source-traceable spectral substitutes. They are not re-analyses of raw historical shot data.
